\myabstract[1]{\vspace{-1cm}\hrule\vspace{2mm}\textbf{Abstract:} This electronic document is a “live” template. The various components of your paper [title, text, heads, etc.] are already defined on the style sheet, as illustrated by the portions given in this document.
\newline
\newline
\textbf{Keywords:} Keyword1, Keyword2, … , KeywordN \vspace{1mm} \hrule}

\maketitle

\section{INTRODUCTION}
\label{sec:introduction_and_overview}

This template is created using Latex, margins, column widths, line spacing, and type styles are built-in; examples of the type styles are provided throughout this document and are identified in italic type, within parentheses, following the example. Some components, such as multi-leveled equations, graphics, and tables are not prescribed, although the various table text styles are provided. The formatter will need to create these components, incorporating the applicable criteria that follow.

\setlength{\parindent}{5mm}

\label{sec:background}

\subsection{Selecting a Template (Heading 2)}
\label{subSec:cmos_power_dissipation}
First, confirm that you have the correct template for your paper size. This template has been tailored for output on the 21cm X 28cm Paper Size.

\subsection{Maintaining the Integrity of the Specifications}
\label{subSec:metrics}

The template is used to format your paper and style the text. All margins, column widths, line spaces, and text fonts are prescribed; please do not alter them. You may note peculiarities. For example, the head margin in this template measures proportionately more than is customary. This measurement and others are deliberate, using specifications that anticipate your paper as one part of the entire proceedings, and not as an independent document. Please do not revise any of the current designations.




\section{USING THE TEMPLATE}
\label{sec:energy_calculation_methodology}


The following is an example of Table~\ref{tab:evolution_of_test_bed}.  Here is also an example of 
Figure~\ref{Fig:ivy_bridge_Chassis} shows VHDL Synthesis Process.


\begin{table*} 
 \center
 \caption{Evolution of test-bed clusters}
    \begin{tabular}{cccccccccc}
    \toprule
    Label & CPU Type & Frequency & cpu\# & core\# & node\# & RAM & HT? & TB? & CS? \\
    \midrule
    A     & Intel\regtm E6550 & 2.00-2.33GHz  & 1  & 2     & 6     & 4GB   & N  & N  & N\\
    B     & Opteron\tm 2354 & 1.1-2.2GHz   & 2  & 4     & 16    & 16GB  & N  & N  & N\\
    C     & Xeon\tm E5-2640v2 & 1.1-2.0GHz & 2  & 8     & 4     & 32GB  & Y  & Y  & Y\\
    \bottomrule
    \end{tabular}
 \label{tab:evolution_of_test_bed}
\end{table*}





\begin{figure}
  \center
  \includegraphics[width=.47\textwidth]{figures/fig1.png}
  \caption{VHDL Synthesis Process}
  \label{Fig:ivy_bridge_Chassis}
\end{figure}

\section{ PREPARE YOUR PAPER BEFORE STYLING}
Before you begin to format your paper, first write and save the content as a separate text file. Keep your text and graphic files separate until after the text has been formatted and styled. Do not use hard tabs, and limit use of hard returns to only one return at the end of a paragraph. Do not add any kind of pagination anywhere in the paper. Do not number text heads-the template will do that for you.
Finally, complete content and organizational editing before formatting. Please take note of the following items when proofreading spelling and grammar:

\subsection{Abbreviations and Acronyms}
Define abbreviations and acronyms the first time they are used in the text, even after they have been defined in the abstract. Abbreviations such as IEEE, SI, MKS, CGS, sc, dc, and rms do not have to be defined. Do not use abbreviations in the title or heads unless they are unavoidable.

\subsection{Authors and Affiliations}
The template is designed so that author affiliations are not repeated each time for multiple authors of the same affiliation. Please keep your affiliations as succinct as possible (for example, do not differentiate among departments of the same organization). 
\section{Bibliography styles}
Here are two sample references:
\cite{gong2021depth}, \cite{murthy2020blockchain}
\section{Conclusions and Future Work}
\label{sec:conclusions_and_future_work}

In this paper we have presented an extensible CPU power measurement framework that supports our own research but is also generally applicable to the computer engineering community in general for accurate computational power consumption measurements.  

